\subsection*{Question 1 \quad Réponse D}

Un produit de facteurs est nul si et seulement si l'un au moins de ses facteurs est nul :
\[(2x - 6)(x + 4) = 0 \iff 2x - 6 = 0 \ \text{ ou } \ x + 4 = 0 \iff x = 3 \ \text{ ou } \ x = -4.\]
D'où $S = \{-4~;~3\}$.

\subsection*{Question 2 \quad Réponse B}

Le dé est équilibré : on est en situation d'équiprobabilité. Les issues strictement
supérieures à 4 sont 5 et 6, soit 2 issues favorables sur 6 :
\[P = \dfrac{2}{6} = \dfrac{1}{3}.\]

\subsection*{Question 3 \quad Réponse C}

On applique $\dfrac{a^m}{a^n} = a^{m-n}$ avec $3 = 3^1$ :
\[\dfrac{3}{3^5} = 3^{1-5} = 3^{-4}.\]

\subsection*{Question 4 \quad Réponse D}

La droite $\mathrm{D}_4$ passe par les points de coordonnées $(-1 ; 0)$ et $(0 ; 1)$. Son ordonnée à
l'origine est $b = 1$ et son coefficient directeur vaut
$a = \dfrac{1 - 0}{0 - (-1)} = 1$. Son équation réduite est donc $y = x + 1$.

\subsection*{Question 5 \quad Réponse B}

Un point appartient à $\mathcal{C}_f$ lorsque son ordonnée est l'image de son abscisse.
Pour H$(-3~;~28)$ :
\[f(-3) = 2\times(-3)^2 - 3\times(-3) + 1 = 18 + 9 + 1 = 28.\]
Le point H appartient donc à $\mathcal{C}_f$.

\subsection*{Question 6 \quad Réponse B}

On réduit d'abord la parenthèse au même dénominateur :
\[E = \left(\dfrac{1}{2} + \dfrac{1}{3}\right)\times\dfrac{2}{5}
= \left(\dfrac{3}{6} + \dfrac{2}{6}\right)\times\dfrac{2}{5}
= \dfrac{5}{6}\times\dfrac{2}{5} = \dfrac{10}{30} = \dfrac{1}{3}.\]

\subsection*{Question 7 \quad Réponse C}

Une hausse de 10\,\% correspond au coefficient multiplicateur $1,10$ et une baisse de
20\,\% au coefficient $0,80$. Le coefficient multiplicateur global est le produit
des deux :
\[1,10 \times 0,80 = 0,88 = 1 - 0,12.\]
Le prix a subi une baisse de 12\,\%.

\subsection*{Question 8 \quad Réponse C}

Toutes les grandeurs sont strictement positives :
\[g = \dfrac{Gm}{r^2} \iff r^2 g = Gm \iff r^2 = \dfrac{Gm}{g} \iff r = \sqrt{\dfrac{Gm}{g}}.\]
